% ============================================================
%  07_pruebas.tex
% ============================================================
\chapter{Pruebas y Resultados}

% ── 7.1 ──────────────────────────────────────────────────
\section{Plan de Pruebas}

Se efectuaron evaluaciones funcionales, de usabilidad y de integración con el
propósito de verificar el cumplimiento de los requisitos, analizar la facilidad de
uso y garantizar la estabilidad del sistema en condiciones reales de operación.

\subsection{Tipos de Pruebas Aplicadas}

\begin{itemize}[leftmargin=1.5em]
  \item \textbf{Pruebas funcionales:} Verificación de que cada módulo responde
        conforme a los requisitos funcionales (RF-01 a RF-12).
  \item \textbf{Pruebas de navegación:} Análisis de la coherencia de los flujos
        de pantalla, menús y transiciones entre roles.
  \item \textbf{Pruebas de integración:} Validación de la comunicación entre
        el cliente React Native y los servicios de Supabase (REST, RPC,
        Realtime, Storage).
  \item \textbf{Pruebas Realtime:} Verificación de la latencia de notificaciones
        en el canal \texttt{orders} y \texttt{table\_calls} bajo condiciones de
        red móvil.
  \item \textbf{Pruebas de seguridad / RLS:} Intentos de acceso cruzado entre
        roles para validar que las políticas RLS bloqueen correctamente.
  \item \textbf{Pruebas exploratorias:} Recorrido libre de la aplicación por
        evaluadores externos para detectar comportamientos inesperados.
  \item \textbf{Pruebas de UX/UI:} Evaluación de la experiencia del usuario,
        priorizando interfaces simples y claras.
\end{itemize}

\subsection{Escenarios Evaluados}

\begin{enumerate}[leftmargin=1.8em]
  \item Registro y login de mesero y administrador.
  \item Creación de un pedido con selección de platillos, mesa y método de pago.
  \item Actualización del estado de un pedido (\texttt{pending} → \texttt{ready}
        → \texttt{delivered}) y verificación de notificación en tiempo real.
  \item Vista de estado de mesas: validar que los cambios de estado aparecen
        sin recargar la pantalla.
  \item Llamada al mesero desde la sesión de cliente en mesa y recepción de
        notificación en el dispositivo del mesero.
  \item Acceso a la carta digital pública (sin login) y validación de que solo
        aparecen ítems con \texttt{available = true}.
  \item Gestión del menú por el admin: crear, editar precio, cambiar disponibilidad,
        subir imagen.
  \item Registro de gastos operativos y generación de reporte financiero por rango
        de fechas.
  \item Consulta del ranking de platos más vendidos (RPC \texttt{get\_top\_items}).
  \item Confirmación de pago QR por el admin y verificación de contabilización
        en el reporte.
  \item Intento de acceso cruzado: mesero intentando acceder a reportes
        (debe ser bloqueado por RLS).
  \item Simulación de pérdida de conexión y reconexión automática del canal
        Realtime.
\end{enumerate}

\subsection{Herramientas de Soporte}

\begin{itemize}[leftmargin=1.5em]
  \item \textbf{Emuladores Android (Android Studio):} Pruebas preliminares de
        compatibilidad y funcionalidad básica.
  \item \textbf{Dispositivos físicos Android:} Validaciones de rendimiento y
        experiencia real del usuario final.
  \item \textbf{Expo Go / EAS Build:} Distribución de builds de prueba al equipo.
  \item \textbf{Supabase Dashboard:} Inspección directa de la base de datos y
        logs de RLS para verificar políticas.
  \item \textbf{Postman:} Pruebas directas de funciones RPC y endpoints REST.
\end{itemize}

\subsection{Condiciones de Aprobación}

Una prueba se considera superada cuando:
\begin{itemize}[leftmargin=1.5em]
  \item Las funcionalidades operan conforme a los requisitos funcionales
        y criterios de aceptación definidos.
  \item Las políticas RLS bloquean accesos no autorizados sin excepción.
  \item Las notificaciones Realtime llegan en menos de 2 segundos en red WiFi.
  \item La experiencia del usuario es valorada como satisfactoria.
  \item No se detectan fallos críticos en las pruebas exploratorias.
\end{itemize}

% ── 7.2 ──────────────────────────────────────────────────
\section{Resultados Obtenidos}

\subsection{Registro y Gestión de Pedidos}

El flujo de creación de pedidos fue validado exitosamente en dispositivos físicos.
El mesero puede seleccionar platillos del menú, asignar una mesa y elegir el método
de pago. El cálculo del total se realiza automáticamente al agregar ítems.

\subsection{Vista de Mesas en Tiempo Real}

La pantalla de estado de mesas del mesero actualizó correctamente los estados
(\texttt{pending}, \texttt{ready}, \texttt{delivered}) al recibir eventos del canal
\texttt{orders} de Supabase Realtime, sin necesidad de recargar la pantalla. La
latencia promedio observada fue inferior a 1.5 segundos en red WiFi.

\subsection{Carta Digital y Sesión de Cliente}

La carta digital pública fue accedida correctamente desde navegadores móviles
sin descarga de app. La sesión anónima por mesa se inició correctamente al
escanear el QR de mesa y el cliente pudo llamar al mesero, recibiendo el mesero
la notificación en tiempo real vía canal \texttt{table\_calls}.

\subsection{Reportes Financieros}

Las funciones RPC \texttt{get\_report} y \texttt{get\_top\_items} devolvieron
resultados correctos en todas las combinaciones de rango de fechas probadas.
La validación de rol dentro de la función \texttt{get\_report} bloqueó
correctamente el acceso cuando se intentó invocar con un token de rol
\texttt{waiter}.

\subsection{Seguridad RLS}

Todos los intentos de acceso cruzado entre roles fueron bloqueados por las
políticas RLS. El mesero no pudo leer órdenes de otros meseros ni acceder
a la tabla \texttt{expenses}. El cliente anónimo solo pudo insertar en
\texttt{table\_calls} para su propia mesa.

% ── 7.3 ──────────────────────────────────────────────────
\section{Limitaciones Conocidas}

\subsection{Free Tier de Supabase}

\begin{center}
\begin{tabular}{ll}
  \toprule
  \textbf{Límite} & \textbf{Valor} \\
  \midrule
  Base de datos                      & 500 MB \\
  Storage                            & 1 GB \\
  Requests mensuales                 & 50,000 \\
  Conexiones Realtime simultáneas    & 200 \\
  Inactividad (pausa automática)     & 7 días sin actividad \\
  \bottomrule
\end{tabular}
\end{center}

\begin{alertbox}
La pausa automática por inactividad es un riesgo operacional real en producción.
Debe resolverse mediante un ping periódico (ej. \texttt{cron-job.org}) o migración
al tier Pro (\$25/mes) antes del primer despliegue en un restaurante real.
\end{alertbox}

\subsection{Limitaciones de Alcance Documentadas}

\begin{itemize}[leftmargin=1.5em]
  \item La \textbf{reasignación de órdenes} entre meseros no está soportada
        en tiempo real. Si el admin reasigna una orden, el nuevo mesero no
        recibe el evento Realtime hasta el próximo ciclo de polling.
  \item El \textbf{cliente en mesa no puede crear órdenes directamente}.
        Requeriría un estado \texttt{draft} y un flujo de confirmación por
        el mesero, fuera del alcance actual.
  \item El \textbf{pago QR es estático} y su verificación es manual por el
        admin. No hay integración con pasarelas de pago en este alcance.
  \item Las \textbf{cancelaciones parciales y devoluciones} no tienen flujo
        dedicado. Se recomienda implementar en versiones posteriores.
\end{itemize}

% ── 7.4 ──────────────────────────────────────────────────
\section{Mejoras Propuestas}

\begin{enumerate}[leftmargin=1.8em]
  \item \textbf{Pasarela de pagos:} Integrar con una pasarela (ej. Mercado Pago
        o Stripe) para verificar pagos QR de forma automática.
  \item \textbf{Dashboard live global:} Implementar una vista de operación en
        tiempo real para el admin con estado de todas las mesas y meseros
        (descartada del alcance actual por complejidad).
  \item \textbf{Cancelaciones parciales:} Añadir flujo para devolver o cancelar
        ítems individuales de una orden sin cancelar toda la comanda.
  \item \textbf{Reasignación de órdenes en tiempo real:} Soporte completo de
        Realtime al reasignar una orden a otro mesero.
  \item \textbf{Soporte multi-idioma:} Internacionalización de la interfaz
        para restaurantes con clientes internacionales.
  \item \textbf{Modo offline básico:} Caché local de las órdenes activas para
        operar sin conexión y sincronizar al reconectar.
\end{enumerate}